\documentclass[11pt]{article}

\usepackage[margin=1in]{geometry}
\usepackage{amsmath,amssymb,amsthm}
\usepackage{enumitem}
\usepackage{hyperref}
\usepackage{xcolor}
\usepackage{titlesec}
\usepackage{booktabs}
\usepackage{array}

\hypersetup{
    colorlinks=true,
    linkcolor=blue!50!black,
    citecolor=blue!50!black,
    urlcolor=blue!50!black
}

\newtheorem{prediction}{Prediction}
\newtheorem{definitionbox}{Definition}

\title{\Large \textbf{A Theory of Computational Coupling Between Intelligent Systems}\\[4pt]
\large Toward a General Foundation for Brain-to-Brain Communication}

\author{\textbf{Ashok Pasala} \\[2pt] \small VIT-AP University \\[2pt] \small \texttt{Independent Research --- Working Draft v0.2.0}}

\date{July 23, 2026}

\begin{document}

\maketitle

\begin{center}
\colorbox{yellow!20}{\parbox{0.9\textwidth}{\small \textbf{Status:} Working Draft (Version 0.2.0). Formal Core (Sections 3--4) completed; empirical evaluation roadmap and foundational literature canon fully specified.}}
\end{center}

\vspace{8pt}

\begin{abstract}
\noindent
Efforts toward brain-to-brain communication (BBI) and related brain-computer interface (BCI) paradigms have largely proceeded by designing communication \emph{protocols} --- fixed encodings, tokenizations, or aligned representational spaces --- without first establishing what quantity such a protocol should maximize, or how to measure whether it has succeeded. We argue this mirrors the state of telecommunications engineering prior to Shannon's definition of channel capacity: a field of designed artifacts without a rigorous target. We introduce a general theory of \textbf{computational coupling} between intelligent systems --- biological or artificial --- that defines a substrate-independent, directed, information-theoretic quantity (\textbf{coupling capacity}) governing how predictively entangled two systems' internal state trajectories can become, given a bandwidth- and structure-constrained interface between them. We derive three falsifiable predictions from this theory, propose estimators suitable for both simulated multi-agent systems and biological recordings, and outline how brain-to-brain communication, communication protocols, and language itself are best understood as special-case applications of this more general quantity, rather than as the object of study in themselves.
\end{abstract}

\section{Introduction}

Prior to 1948, telecommunications engineering possessed extensive practical technique but no rigorous mathematical definition of \emph{information} or of a channel's fundamental capacity. Claude Shannon's landmark contribution \cite{shannon1948mathematical} was not a specific telegraph device or radio modem, but a \emph{measurement theory} that specified the theoretical upper bound of error-free transmission over any noisy medium.

We argue that current efforts in Brain-to-Brain Communication (BBI) and Brain-Computer Interfaces (BCIs) occupy an analogous position today. Substantial engineering ingenuity is devoted to designing tokenization schemes, neural decoders, and aligned embedding spaces, yet the field lacks a governing quantity that specifies:
\begin{enumerate}
    \item What fundamental quantity is being optimized across a brain-to-brain link?
    \item What are the theoretical upper bounds on information exchange between two dynamic, non-stationary neural manifolds?
\end{enumerate}

\subsection{The Core Reframe}
Communication between two intelligent systems (biological or artificial) is fundamentally not the transmission of a static pre-formed message, but the degree to which their internal state trajectories become \textbf{mutually predictive} through a designed or evolved interface map. This reframing subsumes decoding-based and representational-alignment approaches as downstream special cases.

\subsection{Central Object}
Instead of inventing arbitrary brain dictionaries, we introduce \textbf{coupling capacity}: a directed, dynamic generalization of Shannon channel capacity, defined over continuous state trajectories rather than static discrete messages.

\subsection{Scope Statement}
This paper presents a measurement theory, not a system proposal. It does not propose a specific BBI hardware architecture or tokenization scheme; it defines the mathematical target against which all such systems must be evaluated.

\section{Related Work}

The theory of computational coupling synthesizes literature across several key domains:

\begin{itemize}[leftmargin=1.2em]
    \item \textbf{Information-Theoretic Causality \& Transfer Entropy:} Schreiber \cite{schreiber2000measuring} introduced transfer entropy to quantify asymmetric, model-free information transfer between dynamic time series. Massey \cite{massey1990causality} formalized directed information for systems with feedback, laying the foundation for modern measures like directed phase transfer entropy (dPTE). Kraskov et al. \cite{kraskov2004estimating} developed non-parametric $k$-NN mutual information estimators.
    \item \textbf{Brain-to-Brain Interfaces (BBIs) \& Brainets:} Early empirical demonstrations by Pais-Vieira et al. \cite{pais2013brain} (rat-to-rat ICMS) and Rao et al. \cite{rao2014direct} (human EEG-to-TMS motor transfer) proved the physical feasibility of direct brain links. Jiang et al. \cite{jiang2019brainnet} expanded this to multi-person networks (BrainNet). However, these implementations operate under extreme bandwidth constraints ($<1$ bit/sec) and lack a formal measurement capacity.
    \item \textbf{Cross-Subject Representational Alignment:} Modern state-of-the-art methods utilize Optimal Transport, such as Fused Unbalanced Gromov-Wasserstein (FUGW) by Thual et al. \cite{thual2022aligning}, to align functional cortical geometries across subjects.
    \item \textbf{Active Inference \& Hermeneutics:} Friston and Frith \cite{friston2015active} modeled communication as mutual free-energy minimization, demonstrating that interacting brains synchronize to minimize joint prediction error under a shared generative model. Hasson et al. \cite{hasson2010speaker} demonstrated speaker-listener neural coupling during natural communication.
    \item \textbf{Emergent Communication in MARL:} Foerster et al. \cite{foerster2016learning} developed Differentiable Inter-Agent Learning (DIAL), demonstrating that deep multi-agent reinforcement learning agents can invent proprietary communication protocols by backpropagating through noisy channels.
    \item \textbf{Large Brain Foundation Models:} Recent advances in self-supervised learning have produced foundational neural codecs, including LaBraM \cite{jiang2024labram} for EEG and BrainLM \cite{caro2024brainlm} for fMRI, establishing universal latent representations of neural state spaces. Scotti et al. \cite{scotti2024mindeye2} introduced MindEye2 for shared-subject visual perception reconstruction.
    \item \textbf{Representation Transfer Threats:} Nakamura et al. \cite{nakamura2024unsupervised} demonstrated zero-shot representation transfer using hyperspherical embeddings. While impressive, their approach performs passive representational mapping; our theory addresses active, closed-loop state trajectory control over bandwidth-limited channels.
\end{itemize}

\section{Formal Framework}
\label{sec:theory}

\subsection{Systems}

Let each system $i \in \{1, \dots, N\}$ (biological or artificial) be characterized by:

\begin{itemize}[leftmargin=1.2em]
    \item A state space $\mathcal{M}_i$ (a manifold, possibly high-dimensional), with internal state trajectory $x_i(t) \in \mathcal{M}_i$ evolving over a time index $t$ (discrete or continuous).
    \item A self-predictive model $P_i$, used by system $i$ to predict its own future state $x_i(t+\Delta)$ from its own past $x_i(\le t)$.
\end{itemize}

We deliberately do not assume $\mathcal{M}_i$ is shared, aligned, or even comparable in dimension across systems --- a direct departure from hyperalignment-based approaches, which presuppose a common representational geometry.

\subsection{Interfaces}

An \textbf{interface} from system $i$ to system $j$ is a (possibly stochastic) map
\[
g_{i \to j} : x_i(\le t) \;\longmapsto\; s_{i \to j}(t),
\]
producing a signal $s_{i \to j}(t)$ that influences the subsequent evolution of $x_j$. Interfaces are constrained by a bandwidth or structural budget --- e.g., a discrete channel with vocabulary size $K$, a continuous channel with bandwidth $B$, or a hardware-constrained stimulation protocol. Let $\mathcal{G}$ denote the admissible set of interfaces under a given constraint budget.

\subsection{Directed Coupling}

We define the directed coupling from $i$ to $j$, at lag $\Delta$ and under interface $g \in \mathcal{G}$, as the transfer entropy (directed information):
\begin{equation}
\mathrm{TE}_{i \to j}(\Delta; g) \;=\; I\big(x_i(t)\, ;\, x_j(t+\Delta) \,\big|\, x_j(\le t)\big),
\end{equation}
the reduction in uncertainty about $j$'s future state given $i$'s state, beyond what $j$'s own past already provides. This generalizes standard transfer entropy to vector- or manifold-valued states with an explicit lag term, since biological systems exhibit genuine transmission delay and representational drift that a lag-free formulation would misattribute to noise.

\subsection{Coupling Capacity}

\begin{definitionbox}[Coupling Capacity]
The \textbf{coupling capacity} from system $i$ to system $j$ is the supremum of directed coupling over admissible interfaces:
\begin{equation}
C(i \to j) \;=\; \sup_{g \,\in\, \mathcal{G}} \; \mathrm{TE}_{i \to j}(\Delta; g).
\end{equation}
\end{definitionbox}

\noindent This is the direct dynamical analogue of Shannon capacity: where Shannon capacity is a supremum over input distributions of mutual information subject to a power or bandwidth constraint, coupling capacity is a supremum over \emph{interface designs} of directed information between two coupled dynamical systems, subject to an analogous constraint.

\subsection{Bidirectional Coupling and Asymmetry}

Define total coupling as $C_{\text{total}} = C(i \to j) + C(j \to i)$, and an asymmetry index
\begin{equation}
A \;=\; \frac{C(i \to j) - C(j \to i)}{C(i \to j) + C(j \to i)} \;\in\; [-1, 1].
\end{equation}
$A$ is predicted to correlate with task-defined role (e.g., ``leading'' vs.\ ``following,'' ``speaking'' vs.\ ``listening''); see Prediction 3 below.

\section{Falsifiable Predictions}
\label{sec:predictions}

\begin{prediction}[Capacity--Bandwidth Law]
For interfaces parameterized by a scalar bandwidth $B$, $C(i \to j; B)$ is monotonically increasing and concave in $B$, saturating at a value set by $\min\big(\dim_{\mathrm{eff}}(\mathcal{M}_i),\, \dim_{\mathrm{eff}}(\mathcal{M}_j)\big)$ --- the smaller of the two systems' \emph{effective} representational dimensionality --- rather than by the channel's raw capacity. This is directly falsifiable: if measured capacity continues scaling with channel bandwidth well past the point where either system's effective dimensionality should cap it, the theory as stated is wrong.
\end{prediction}

\begin{prediction}[Self-Predictive Accuracy Governs Capacity Efficiency]
Let $R_i$ denote system $i$'s self-predictive accuracy --- its own model's log-likelihood in predicting its own future state. Capacity achieved per unit bandwidth, $C(i \to j)/B$, is monotonically increasing in $R_i$ and $R_j$ jointly: systems with better internal world models extract more coupling from an identical channel.
\end{prediction}

\begin{prediction}[Asymmetry Tracks Role]
The asymmetry index $A$ correlates with an externally defined task-role variable across domains (multi-agent reinforcement learning, human dialogue, hyperscanning) --- a cross-domain claim, not a redescription of a single-domain result.
\end{prediction}

\section{Evaluation Metrics and Estimation}

Estimating coupling capacity in continuous high-dimensional state spaces requires estimators tailored to the data regime:

\begin{itemize}[leftmargin=1.2em]
    \item \textbf{Simulated Multi-Agent Systems (Ground Truth):} In artificial multi-agent reinforcement learning (MARL) testbeds, full access to internal recurrent states $h_t^A, h_t^B$ permits direct estimation via $k$-nearest-neighbor ($k$-NN) Kraskov-St\"ogbauer-Grassberger (KSG) transfer entropy estimators. Alternatively, we introduce a \emph{predictive-gain estimator}:
    $$\hat{\mathrm{TE}}_{A \to B} = L_{\text{self}} - L_{\text{joint}} = \mathbb{E} \left[ \log \frac{P_\psi(x_B(t+1) \mid x_B(\le t), x_A(\le t))}{P_\phi(x_B(t+1) \mid x_B(\le t))} \right]$$
    which scales effectively to deep neural networks without suffering from the curse of dimensionality.
    \item \textbf{Biological Recordings (EEG / fMRI Hyperscanning):} State trajectories are extracted via non-linear manifold learning or pre-trained foundation model encoders (LaBraM, BrainLM) before applying the predictive-gain estimator. Effective Transfer Entropy (ETE) is computed by subtracting surrogate-shuffled baseline estimates to eliminate finite-sample bias.
\end{itemize}

\section{Falsifiability, Scope, and Discussion}

\subsection{Falsifiability \& Scope}
A critical requirement for any measurement theory is explicit commitment to falsifiability:
\begin{itemize}[leftmargin=1.2em]
    \item \textbf{Falsification Criteria:} If measured transfer entropy across an optimized channel scales indefinitely with raw channel bandwidth $B$ past the effective representational dimensionality cap $\min(\dim_{\text{eff}}(\mathcal{M}_i), \dim_{\text{eff}}(\mathcal{M}_j))$, Prediction 1 is falsified, and the theory as stated is invalid.
    \item \textbf{Scope Limitation:} This framework quantifies \emph{how much} two dynamic systems become predictively entangled, not \emph{what} semantic content is transmitted. Content-level interpretation is a downstream task dependent on specific representational geometries.
\end{itemize}

\subsection{Limitations \& Ethical Considerations}
\begin{itemize}[leftmargin=1.2em]
    \item \textbf{Estimator Sample Complexity:} High-dimensional transfer entropy estimation remains computationally intensive ($O(N \log N)$ for $k$-NN estimators). Neural predictive-gain models mitigate this but introduce optimization hyperparameter dependencies.
    \item \textbf{Non-Stationarity of Neural Manifolds:} Biological neural manifolds exhibit non-stationarity over extended timeframes due to plasticity, requiring adaptive windowing in transfer entropy estimators.
    \item \textbf{Ethical Considerations for Stimulation Work:} Eventual bidirectional closed-loop BBI prototypes involving non-invasive TMS or invasive microstimulation introduce ethical questions regarding agency, cognitive privacy, and involuntary mental state alignment.
    \item \textbf{Comparison with Integrated Information Theory ($\Phi$):} While Giulio Tononi's IIT quantifies integrated information \emph{within} a single system, Coupling Capacity $C_{\text{couple}}$ quantifies directed information flow \emph{between} dynamic systems across a bandwidth-limited interface.
\end{itemize}

\section{Multi-Year Empirical Roadmap}

\begin{enumerate}[leftmargin=1.4em]
    \item \textbf{Paper 1 (Multi-Agent RL):} Validate Predictions 1 and 2 in controlled cooperative games (\texttt{PettingZoo} \texttt{simple\_speaker\_listener\_v4}), sweeping channel quantization $B \in [1, 2, 4, 8, 16, 32]$ bits/step.
    \item \textbf{Paper 2 (Hyperscanning):} Test Predictions 1--3 against public human dyadic EEG corpora (OpenNeuro \texttt{ds007764} DUET, \texttt{ds007471} Joint Agency EEG) using FUGW functional alignment.
    \item \textbf{Paper 3 (Causal Manipulation):} Causal non-invasive intervention manipulating feedback latency and bandwidth in human dyads.
    \item \textbf{Paper 4 (Learned BCP Prototype):} Use $C_{\text{couple}}$ as an explicit differentiable loss function to train a closed-loop BBI prototype.
\end{enumerate}

\section{Conclusion}

We have introduced a general Theory of Computational Coupling between intelligent systems, reframing Brain-to-Brain Communication (BBI), communication protocols, and language from designed engineering artifacts into testable applications of a single underlying quantity: \textbf{Coupling Capacity} ($C_{\text{couple}}$). By generalizing Shannon's channel capacity to dynamic, self-predictive state trajectories on neural manifolds, this framework provides the field with a rigorous measurement target.

\vspace{12pt}
\bibliographystyle{plain}
\bibliography{references}

\end{document}
